\documentclass[11pt]{article}
\usepackage[margin=1in]{geometry}
\usepackage[T1]{fontenc}
\usepackage{xcolor}
\usepackage{enumitem}
\usepackage{booktabs}
\usepackage{tabularx}
\usepackage[colorlinks=true,linkcolor=blue,urlcolor=blue]{hyperref}
\setlength{\parindent}{0pt}
\setlength{\parskip}{6pt}
\definecolor{revisionyellow}{RGB}{255,244,153}
\newcommand{\reviewercomment}[1]{%
  \vspace{4pt}\noindent\colorbox{revisionyellow}{%
  \parbox{\dimexpr\linewidth-2\fboxsep}{\textbf{Reviewer comment:} #1}}}
\newcommand{\response}{\noindent\textbf{Response:} }

\begin{document}
\begin{center}
{\Large\bfseries Response to Reviewer}\\[6pt]
{\bfseries Manuscript: ``Anomaly Detection Algorithm for Enhanced Fraud
Identification in the Insurance Sector''}\\[2pt]
{\itshape Annals of Data Science --- Major Revision}
\end{center}

Dear Editor and Reviewer,

We sincerely thank you for the careful and constructive assessment. The
review helped us sharpen the manuscript's conceptual contribution, replace the
chronological literature list with a systematic taxonomy, and present the
existing experiments as an explicit component and threshold ablation. All
new or substantially rewritten passages are identified by yellow revision
markers in the highlighted manuscript; a clean manuscript is also provided.
No numerical result has been invented or inferred from unavailable
prediction-level data.

\reviewercomment{The manuscript must more sharply delineate the core
innovations of SAD. The novelty of its ``sequential'' and ``multilayer''
structures is insufficiently highlighted; a comparative table versus
state-of-the-art counterparts is recommended.}

\response We have substantially revised the Abstract, Introduction, Literature
Review, Methodology, and Conclusion. The revision now defines
\emph{sequential} as partition-wise weak-to-strong supervision: anomaly
consensus produces the first candidate set, verified labels initialize the
supervised ensemble, and the ensemble is retrained after every later
partition. It defines \emph{multilayer} as an ordered cascade: anomaly
consensus first constructs a shortlist, EAD then acts as a second-stage gate,
and only the surviving claims enter verification. We also state explicitly
that SAD is neither a temporal sequence model nor a reinforcement-learning
agent because it contains no state--action--reward formulation.

The Introduction now contains a four-part contribution statement. In addition,
new Table~1 (``Positioning of SAD relative to representative fraud-detection
approaches'') compares SAD with supervised/cost-sensitive, unsupervised,
reinforcement/active-verification, and hybrid/graph approaches in terms of
learning signal, principal strength, and the limitation targeted by SAD. The
formal threshold notation and corrected pseudocode make the stage ordering
unambiguous.

\reviewercomment{The literature review needs restructuring. The list-like
survey should be reorganized into supervised, unsupervised,
reinforcement-learning-based, and hybrid architectures, with the strengths,
limitations, and research gaps of each class summarized explicitly.}

\response We have completely replaced the chronological survey with the four
requested categories:
\begin{enumerate}[label=(\roman*),leftmargin=*]
  \item supervised approaches;
  \item unsupervised approaches;
  \item reinforcement- and active-learning-based approaches; and
  \item hybrid and graph-based architectures.
\end{enumerate}
Each subsection now explains the corresponding learning signal, practical
strengths, limitations, and relationship to costly fraud verification. A new
``Comparative positioning and research gap'' subsection synthesizes the gap:
an end-to-end process that begins with scarce labels, constrains verification
workload, converts verified claims into supervised training data, and updates
the ensemble as batches arrive. We also added and verified recent
insurance-specific references from 2024--2026, including systematic-review,
class-imbalance, interpretable boosting, graph-feature, and deep-learning work.

\reviewercomment{The experimental evaluation is incomplete. Recent
fraud-detection-specific baselines are missing, and no ablation study or
validation on an independent second dataset is provided; adding these
experiments would strengthen generalizability.}

\response We agree that these additions are important and have addressed this
comment to the fullest extent supported by the submitted research artifacts.

First, the revised literature taxonomy and comparison table add recent,
fraud-specific reference points (including Schrijver et al., 2024;
Banulescu-Radu and Yankol-Schalck, 2024; Khalil et al., 2024; Vorobyev, 2024;
Ding et al., 2025; and Yankol-Schalck, 2026). The empirical table already
compares SAD with logistic regression, SVM, XGBoost, KNN, random forest,
AdaBoost, naive Bayes, and an artificial neural network. We now distinguish
these implemented baselines from recent literature comparators and do not
claim numerical head-to-head results that were not run under the same data
split.

Second, we reorganized the existing architecture and threshold experiments as
a formal component ablation. Simple-A omits EAD; Simple-B inserts EAD into the
flat anomaly ensemble; and Multilayer places EAD after consensus as a separate
gate. All three retain the same supervised-model set, partition schedule, and
normal-sampling percentage. The revised interpretation reports concrete
contrasts, including the accuracy difference at
$(T_{ML},T_{AD})=(4,3)$ and the first-partition runtime reduction from 25.18
minutes (Simple-B) to 1.78 minutes (Multilayer). We also recast the threshold
grid as sensitivity analysis and explicitly avoid claiming monotonicity where
the table does not show it.

Third, the available source package contains only the Carclaims benchmark and
does not include prediction-level outputs or access to a genuinely independent
claims portfolio. A differently named copy of the same benchmark would not be
independent validation. Rather than overstate generalizability, we have
restricted every empirical claim to within-dataset evidence and added an
``External validity and remaining experimental needs'' subsection. It specifies
the proper next protocol: freeze preprocessing and thresholds on Carclaims,
test them on a genuinely independent insurer/product/time-period dataset, and
report repeated-seed intervals, fraud-class confusion matrices, PR-AUC, review
rate, and investigation cost. We respectfully present this limitation
transparently because manufacturing a second-dataset result would be
scientifically inappropriate.

\textbf{Additional corrections.} During revision we also corrected the sign of
the Shannon-entropy definition, restricted SMOTE to training data, corrected
the row-level scope of the pseudocode, replaced ``ML'' with ``LR'' in the
results table, aligned figure captions with the stated partitions, clarified
that the recorded target simulates expert feedback, and removed unsupported
reinforcement-learning and deployment claims.

We hope that these revisions make the novelty, evidence, and present limits of
SAD substantially clearer. Thank you again for helping us improve the
manuscript.

Sincerely,\\
Mehrab Atighi, Sayed Jamal Mirkamali, and Amir T. Payandeh Najafabadi

\end{document}
